\section{Character-valued Stieltjes moments and Euler-factor cancellations}
\label{stchar:section}

The character-theoretic question in the colored-moment report
\cite{stcol:report} asks for an exact Stieltjes dictionary at the trivial
zeros of Dirichlet $L$-functions, including imprimitive characters.
The balanced scalar transform gives such a dictionary.  Primitive
characters lead to positive-order $L$-values; omitted Euler factors
produce additional order-zero identities and a family of vanishing
ordinary moments.  The functional equation and Euler products used
here are classical \cite{stchar:classical}.

\subsection{The character kernel and its complete moment criterion}

Let $\chi$ be a nonprincipal Dirichlet character modulo $d$, and let
$\chi_*$ be its primitive inducing character of conductor $f>1$.
For $q\ge0$ define
\begin{align}
 K_{q,\chi}(x)&=\frac1d\sum_{r=1}^d\chi(r)
              C_q\left(\frac xd,\frac rd\right),\qquad x>0,
              \label{stchar:kernel}\\
 &=\frac1d\sum_{r=1}^d\chi(r)\sum_{j=0}^q(-1)^j\binom qj
       \left(\log\frac xd\right)^{q-j}
       \gamma_j\left(\frac{x+r}{d}\right).\notag
\end{align}
The equality uses $\sum_r\chi(r)=0$ to cancel the extra uncolored
logarithmic term.  Thus every formula below involves a finite, explicit
combination of generalized Stieltjes functions.  Equivalently, put
\[
 F_\chi(s,x)=\sum_{n\ge1}\frac{\chi(n)}{(x+n)^s}
    =d^{-s}\sum_{r=1}^d\chi(r)\zeta\left(s,\frac{x+r}{d}\right).
\]
The series is initially taken for $\Re s>1$, and the finite Hurwitz
expression continues it to an entire function of $s$.  Then
$K_{q,\chi}(x)=\partial_z^q[x^zF_\chi(1+z,x)]_{z=0}$.

\begin{theorem}[Character Mellin transform and ordinary moments]
\label{stchar:mellin}
For $0<\Re u<q+1$,
\begin{equation}\label{stchar:transform}
 \int_0^\infty x^{u-1}K_{q,\chi}(x)\,dx
   =\frac{\pi}{\sin\pi u}Q_q(u)L(1-u,\chi).
\end{equation}
For an integer $m\ge q+1$, the ordinary integral on the left
converges if and only if
\begin{equation}\label{stchar:zero-block}
 L(-q,\chi)=L(-q-1,\chi)=\cdots=L(1-m,\chi)=0.
\end{equation}
It then converges absolutely and equals
\begin{equation}\label{stchar:moment}
 (-1)^{m+1}Q_q(m)L'(1-m,\chi).
\end{equation}
\end{theorem}
\begin{proof}
The periodic density
\[
 \rho_\chi(t)=\sum_{n\ge1}\chi(n)e^{-nt}
       =\frac{\sum_{r=1}^d\chi(r)e^{-rt}}{1-e^{-dt}}
\]
is analytic at zero, because the numerator vanishes there, and
decays exponentially at infinity.  Its Taylor coefficients are
$(-1)^kL(-k,\chi)/k!$, as follows by expanding the finite Hurwitz
formula, or its Bernoulli generating function.  The proof of
Lemma~\ref{stcol:tail-lemma} therefore gives
\begin{equation}\label{stchar:tail}
 K_{q,\chi}(x)
  =\sum_{k=q}^N\frac{(-1)^kL(-k,\chi)}{k!}
       \left.\partial_z^q(1+z)_k\right|_{z=0}x^{-k-1}
       +O(x^{-N-2}).
\end{equation}
The finite Stieltjes sum is $O(1+|\log x|^q)$ at zero.
For $0<\Re u<1$, absolute Fubini in the density integral yields
$\Gamma(u+z)\Gamma(1-u)L(1-u,\chi)/\Gamma(1+z)$.
Differentiate and continue to the strip permitted by
\eqref{stchar:tail}.  This proves \eqref{stchar:transform}.
The nonzero Pochhammer coefficients prove the necessary and sufficient
zero-block condition exactly as in Theorem~\ref{stcol:exceptional}.
Taking the removable value at $u=m$ gives \eqref{stchar:moment}.
\end{proof}

\subsection{Primitive characters: evaluated identities at every trivial zero}

Write $\chi_*(-1)=(-1)^\varepsilon$, $\varepsilon\in\{0,1\}$,
and use the Gauss sum convention
\[
 \tau(\chi_*)=\sum_{r=1}^f\chi_*(r)e^{2\pi ir/f}.
\]
For $q=\varepsilon+2n$ with $n\ge0$, define
\begin{equation}\label{stchar:constant}
 T_q(\chi_*)=
   \frac{(-1)^n i^\varepsilon\tau(\chi_*)q!f^q}
        {2(2\pi)^q}L(q+1,\overline{\chi_*}).
\end{equation}
\begin{theorem}[Primitive exceptional-moment dictionary]
\label{stchar:primitive}
For a primitive nonprincipal character and $q=\varepsilon+2n$,
\begin{equation}\label{stchar:primitive-moment}
 \boxed{\displaystyle\int_0^\infty x^qK_{q,\chi_*}(x)\,dx
       =T_q(\chi_*).}
\end{equation}
These are all its extra ordinary integer moments: precisely one exists
at orders $q\equiv\varepsilon\pmod2$, and none at the other orders.
\end{theorem}
\begin{proof}
The completed functional equation is
\[
 \Lambda(s,\chi_*)=
  \left(\frac f\pi\right)^{(s+\varepsilon)/2}
       \Gamma\left(\frac{s+\varepsilon}2\right)L(s,\chi_*),
 \qquad
 \Lambda(s,\chi_*)=
       \frac{i^{-\varepsilon}\tau(\chi_*)}{\sqrt f}
       \Lambda(1-s,\overline{\chi_*}).
\]
At $s=-q=-\varepsilon-2n$, the gamma factor has residue
$2(-1)^n/n!$.  Equating the finite limits gives
\[
 L'(-q,\chi_*)=
 \frac{(-1)^n i^{-\varepsilon}\tau(\chi_*)q!f^q}
      {2(2\pi)^q}L(q+1,\overline{\chi_*}).
\]
Here $n!\Gamma(n+\varepsilon+\tfrac12)=q!\sqrt\pi/2^q$.
Since $Q_q(q+1)=1$, \eqref{stchar:moment} multiplies this derivative
by $(-1)^q=(-1)^\varepsilon$, proving the formula.
The functional equation and nonvanishing of $L(s,\overline{\chi_*})$
for real $s>1$ show that the negative-integer zeros are precisely
the indicated simple trivial zeros.  At $q=0$, the same assertion
uses the classical nonvanishing $L(1,\overline{\chi_*})\ne0$.
Consecutive negative-integer values therefore cannot both vanish.
\end{proof}

For example, let $\chi_3(1)=1$, $\chi_3(2)=-1$ modulo three,
and let $\chi_5$ be the quadratic character modulo five.  Their
Gauss sums are $i\sqrt3$ and $\sqrt5$, respectively.  Formula
\eqref{stchar:primitive-moment} gives
\begin{align}
 \int_0^\infty xK_{1,\chi_3}(x)\,dx
     &=-\frac{3\sqrt3}{4\pi}L(2,\chi_3),\label{stchar:mod3}\\
 \int_0^\infty x^2K_{2,\chi_5}(x)\,dx
     &=-\frac{25\sqrt5}{4\pi^2}L(3,\chi_5).
       \label{stchar:mod5}
\end{align}

\subsection{Imprimitive characters and the exceptional order-zero block}

Put
\begin{equation}\label{stchar:euler-factor}
 \mathcal P=\{p:p\mid d,\ p\nmid f\},\qquad
 E_\chi(s)=\prod_{p\in\mathcal P}(1-\chi_*(p)p^{-s}),\qquad
 L(s,\chi)=E_\chi(s)L(s,\chi_*).
\end{equation}
Repeated prime powers in the modulus have no additional effect.
The exact kernel counterpart is the dilation identity
\begin{equation}\label{stchar:dilation}
 K_{q,\chi}(x)=
 \sum_{h\mid\prod_{p\in\mathcal P}p}
       \frac{\mu(h)\chi_*(h)}h K_{q,\chi_*}(x/h).
\end{equation}
Indeed inclusion--exclusion in the defining Dirichlet series writes
\[
 F_\chi(s,x)=\sum_h\mu(h)\chi_*(h)h^{-s}F_{\chi_*}(s,x/h).
\]
The factor $x^zh^{-z}$ becomes $(x/h)^z$, so differentiation
does not introduce extra logarithmic terms in \eqref{stchar:dilation}.

For every $q\ge1$, the factors $1-\chi_*(p)p^q$ are nonzero.
Thus an extra ordinary moment exists precisely for
$q\equiv\varepsilon\pmod2$, and its value is
\begin{equation}\label{stchar:imprimitive-moment}
 \int_0^\infty x^qK_{q,\chi}(x)\,dx
     =E_\chi(-q)T_q(\chi_*).
\end{equation}
There are no longer zero blocks starting at $q\ge1$.

To give the complete order-zero statement, let
\[
 r=\#\{p\in\mathcal P:\chi_*(p)=1\}.
\]
If $\varepsilon=0$, then $L(0,\chi)=0$ and $L(-1,\chi)\ne0$.
There is precisely one extra moment, at $m=1$, with value
$E_\chi(0)T_0(\chi_*)$, possibly zero.  If $\varepsilon=1$ and
$r=0$, there is no extra order-zero moment.  If $\varepsilon=1$
and $r\ge1$, then
\begin{equation}\label{stchar:odd-block}
 L(0,\chi)=L(-1,\chi)=0,\qquad L(-2,\chi)\ne0.
\end{equation}
There are exactly two extra moments, and
\begin{align}
 \int_0^\infty K_{0,\chi}(x)\,dx
    &=E_\chi'(0)L(0,\chi_*),\label{stchar:zero-moment}\\
 \int_0^\infty xK_{0,\chi}(x)\,dx
    &=\int_0^\infty xK_{1,\chi}(x)\,dx
      =E_\chi(-1)T_1(\chi_*).\label{stchar:first-moment}
\end{align}
Here $E_\chi'(0)=0$ if $r\ge2$, while if the sole omitted prime
with value one is $p_0$, then
\[
 E_\chi'(0)=\log p_0
       \prod_{p\in\mathcal P\setminus\{p_0\}}(1-\chi_*(p)).
\]
These claims follow because every factor with $\chi_*(p)=1$
has a simple zero at $s=0$.  For odd $\chi_*$ the value
$L(0,\chi_*)$ is nonzero, whereas $L(-1,\chi_*)$ is a simple
trivial zero.  Theorem~\ref{stchar:mellin} then proves every
convergence assertion and value, including the ordinary vanishing identity
\begin{equation}\label{stchar:vanishing}
 \boxed{\displaystyle\int_0^\infty x
       [K_{0,\chi}(x)-K_{1,\chi}(x)]\,dx=0
       \quad(\varepsilon=1,\ r\ge1).}
\end{equation}
No regularized integral is used.  All the integrals in
\eqref{stchar:zero-moment}--\eqref{stchar:vanishing} converge absolutely.

For a concrete digamma identity, set
\[
 H(x)=K_{0,\chi_3}(x)
     =\frac13\left[\psi\left(\frac{x+2}3\right)
                    -\psi\left(\frac{x+1}3\right)\right].
\]
The character induced from $\chi_3$ modulo $21$ removes the prime seven.
Since $\chi_3(7)=1$ and $L(0,\chi_3)=1/3$, the formulas give
\begin{align}
 \int_0^\infty\left[H(x)-\frac17H(x/7)\right]dx
      &=\frac{\log7}3,\label{stchar:21zero}\\
 \int_0^\infty x\left[H(x)-\frac17H(x/7)\right]dx
      &=\frac{9\sqrt3}{2\pi}L(2,\chi_3).
        \label{stchar:21first}
\end{align}
Removing both seven and thirteen, equivalently inducing modulo $273$,
gives the further exact pair
\begin{align}
 \int_0^\infty\left[H(x)-\frac17H(x/7)-\frac1{13}H(x/13)
                         +\frac1{91}H(x/91)\right]dx&=0,
                         \label{stchar:273zero}\\
 \int_0^\infty x\left[H(x)-\frac17H(x/7)-\frac1{13}H(x/13)
                         +\frac1{91}H(x/91)\right]dx
     &=-\frac{54\sqrt3}{\pi}L(2,\chi_3).
                         \label{stchar:273first}
\end{align}
The bracketed kernels decay as $O(x^{-3})$.  Individual unweighted
integrals of $H$ diverge; the formulas concern their displayed combined
kernels.  More generally this character construction answers the
primitive/imprimitive exceptional-moment question for all nonprincipal
characters.  Principal characters, whose inducing function has a pole
at one, require a different scalar subtraction; their completion is
given in Section~\ref{stdil:section}.
